Ear images captured in unconstrained settings are small and JPEG-compressed, whereas efficient super-resolution (SR) networks are designed and ranked on large natural images that are downsampled bicubically and never compressed. We measure what this mismatch costs and where it should be repaired. On a benchmark of ear images with clean ground truth at three input sizes (24 to 48 pixels) and on two in-the-wild ear datasets, the \NPub{} fidelity-oriented SR networks with public weights that we evaluate, including the NTIRE 2026 efficient-SR baseline and five of its entries, lose their advantage over bicubic interpolation once the input is compressed: at JPEG quality 75 all of them fall below bicubic on the controlled benchmark (by \AmiQLowGapMin{} to \AmiQLowGapMax{}~dB), while at quality 93 all of them are above it again. Real small ear images sit on both sides of this threshold. Under compression the \NPub{} architectures differ by only \SpreadQLow{}~dB and a \ParamsRatio{}$\times$ larger network is no better, whereas changing the training degradation of a single backbone changes its PSNR by \DegEffMin{} to \DegEffMax{}~dB. We therefore repair the degradation rather than the architecture: we fine-tune an off-the-shelf real-time network with a compression-aware degradation measured from real small ear images, using a photometric augmentation and in-the-wild ear images without which fine-tuning on laboratory data collapses on bright inputs. On synthesised inputs from three datasets the resulting model exceeds bicubic interpolation by \EstVsBicMin{} to \EstVsBicMax{}~dB, the published weights by \EstVsPubMin{} to \EstVsPubMax{}~dB and a generic real-world degradation pipeline by \EstVsGenMin{} to \EstVsGenMax{}~dB, at an unchanged cost of \LatSpan{}~ms per image on a desktop GPU. On real small ear images that are JPEG-compressed, for which no ground truth exists, rank-1 identification rises from \RecRankBic{}\% with bicubic upscaling and \RecRankPub{}\% with the published weights to \RecRankEst{}\%, and the published weights help only where the files are lightly compressed; on a second dataset whose images carry no compression structure, no upscaling method changes identification. An ablation shows that covering the range of compression levels accounts for the gain; the measured noise and blur do not add to it.
